% !PRINTED 235
\weberpage{235}

%<<CONT
\gesperrt{Determinante} $-1$ \gesperrt{als Substitutionen der zweiten Art} bezeichnen.

Nach den Resultaten des \S.\,50 sind diese Gruppen isomorph mit den Gruppen bin\"arer linearer Substitutionen mit der Determinante $+1$ und $-1$, w\"ahrend bei den gebrochenen Substitutionen, die uns auf die Poly\"edergruppen gef\"uhrt haben, die beiden Arten nicht unterscheidbar sind.

Die endlichen Gruppen der zweiten Art sind, abgesehen von ihrem allgemeinen gruppentheoretischen Interesse, wichtig wegen ihrer Anwendung in der Krystallographie. Ihre vollst\"andige Bestimmung bietet jetzt keine wesentlichen Schwierigkeiten mehr. Wir verstehen unter Substitutionen schlechtweg bin\"are lineare Substitutionen mit der Determinante $\pm 1$, und erinnern daran, dass zwei solche Substitutionen, deren Coëfficienten sich nur durch das gemeinschaftliche Vorzeichen unterscheiden, wie
\[
\begin{pmatrix} \alpha, & \beta \\ \gamma, & \delta \end{pmatrix}, \quad
\begin{pmatrix} -\alpha, & -\beta \\ -\gamma, & -\delta \end{pmatrix},
\]
als nicht verschieden gelten (\S.\,49).

Da eine Gruppe linearer Substitutionen die identische Substitution enth\"alt, die von der Determinante $+1$ ist, so m\"ussen in jeder Gruppe $Q$ der zweiten Art neben den uneigentlichen auch eigentliche Substitutionen vorkommen, und diese eigentlichen Substitutionen bilden f\"ur sich eine Gruppe $P$. Ist $\varphi$ irgend eine in $Q$ vorkommende Substitution der zweiten Art, so kommt jede Composition von $\varphi$ mit einer Substitution aus $Q$ der zweiten Art in $P$ vor, und daher k\"onnen wir $Q$ immer so zerlegen:
\weq{1}{Q = P + P\varphi.}
Die Gruppe $P$ muss eine der fr\"uher betrachteten Poly\"edergruppen sein. Es muss $\varphi^{-1}P\varphi = P$ sein, und $P$ ist also ein \gesperrt{Normaltheiler} von $Q$. Ist umgekehrt $P$ eine Poly\"edergruppe und $\varphi$ eine Substitution zweiter Art, die der Bedingung $\varphi^{-1}P\varphi = P$ gen\"ugt, so ist $Q = P + P\varphi$ eine \gesperrt{Gruppe der zweiten Art}.

Auch hier betrachten wir zwei Gruppen, die durch Transformation aus einander hervorgehen, als nicht wesentlich verschieden. Wir haben, um alle $Q$ zu bilden, die verschiedenen F\"alle durchzugehen.

I. Ist $P$ die Einheitsgruppe, die wir als cyklische Gruppe ersten Grades mit $C_1$ bezeichnen, besteht also $P$ nur aus der
%CONT>>
